\pg{335}  -"essencia: Ora: o y ente, que: hão existe dia: exigeneiar d "sua. est:

\textasciicircum{} Sencia, consta dé- existencia-realniente--- istincta--da----essencia;- e -por isso de \emph{acto} e de \emph{potencia}. Logo o ente, que E \emph{contingente}, \textbar{} consta de \emph{acto} e de \emph{potencia} \footnote{; \& essencia do ente \emph{contingente}, mesmo quando « €' actual e existe -na ' realidade, distingue-se realmente-da- \emph{existencia}, como- a-potencia- eal-se-.di-.stingue do acto real. Portanto o ente \emph{contingente} consta: de \emph{acto} e de \emph{potencia}. --- D'ahi se segue que o ente \emph{contingente} não é necessariamente. unico, mas podem.existir muitos entes contingentes. "Por quanto, : Ó ser: do \emph{contingente}, sendo recebido na essencia, pode multiplicar-se indefinitàmente, conforme a multidão dos sujeitos, ou das essencias, em:que é recebido: \textbar{} -}.. .

58::0 ente necessario. é eterno. \textbar{}que, pela: exigencia..da -sua essencia, não: pode. deixar: de: existir. Ora o ente, que não "pode" deixar: de existir; existia sempre; exi- -

59. 0 ente. contingente € “produzido” por uma i causa extrin-. seca; --- Tudo o: que se encontra num-sujeito;-e- não: “pertence: á sua essencia, deve ser produzido - --- ou pelos principios essenciaes 5 do. proprio sujeito. --- ou: por: uma causa extrinseca: Não h meio termo Ora ; a existencia, que. não per tene :

“essência do ente "\emph{conti gente}, "não pode-ser- produzida \textbf{pelos-principios-essenciáesdo} proprio ete; porque, neste caso; a essencia preexistiria-a ' existencia (pois a causa: ;preexiste. ao. effeito)

que é absurdo. Logo deve ser. produzida. por uma. causa. extrinseca:. Logo. o. ente ' \emph{contingente} é produzido” por uma causa: extrinseca \emph{O}:\footnote{O ente necessario, não podendo 'deixar de ser, é eterno; porque, se não fosse eterno, se tivesse começo ou fim, seria \emph{contingente}. --- Além d'isso, o ente necessario, sendo \emph{acto purissimo}, não pode ter nem principio, nem fim, . nem Successão, porque tudo isto envolve potencialidade; e por isso € \emph{eterno}, porque a eternidade exclue prin\&ipio, fim e successão. --- Não só o ente \emph{necessario} é eterno, mas é a sua \emph{eternidade}; porque elle é o seu ser, e a sua duração ou permanencia no \emph{ser} não se distingue da sua essencia, e por isso, assim como é a sua essencia, assim tambem é a sua eternidade. (Sum. \emph{Th}. p. 1, q. 10, a. 2; c. Gent., 1. I, c. 15).\par (3 O ente \emph{contingente}, por não ser o proprio ser ou a propria existencia, isto é, por não ser \emph{necessario}, € \emph{indifferente} para \emph{o} existencia e para a nàoexistencia\pgn{336}, e por isso não pode ser eterno, 'e teve começo. --- E como recebeu elle a existencia real? Foi pela exigencia da sua natureza? Não. Porquanto, um sujeito. que é por si indifferente para dois attributos oppostos, não pode ter na sua essehcia a razão. pela. qual exija um attributo, e não o opposto. Ora o ente \emph{contingente} é. pela sua, essencia, indifferente para a existencia e para a não-existencia, é por isso a existencia não. lhe pode convir.pela mesma, essencia. Logo o.ente contingente deve receber de um agente extrinseco a existencia real. (\emph{Sum}. \emph{Th}. p. 1, q. 3, a. 4). -}

